\subsection{Interpretation of Main Results}
\label{sec:discussion-interpretation}

The quantitative results in Tables~\ref{tab:denoising-metrics} and~\ref{tab:feature-preservation} reveal an important insight: the performance gap between PMDF-Net and all baseline methods cannot be attributed to any single architectural or loss-function choice alone, but rather to their deliberate combination within a unified multi-domain framework. PMDF-Net's $+11.4$\,dB SNR improvement is $1.3$\,dB higher than the best single-domain baseline, a margin that may appear modest in absolute terms but represents a $>35\%$ reduction in residual noise variance. More significantly, this SNR gain does not come at the expense of waveform fidelity---PMDF-Net simultaneously achieves the lowest arrival time error ($0.14$\,$\mu$s) and the highest envelope concordance ($0.981$) among all methods. This joint optimization of denoising strength and feature preservation is the hallmark of the proposed approach.

The mechanism underlying this dual achievement lies in the complementary information captured by the time and time-frequency branches. The time-domain branch encodes precise temporal dynamics, including arrival times and wave packet morphology, while the time-frequency branch encodes frequency content that changes characteristically when defects alter local wave speed and thickness. A single-domain architecture must compress both information streams into one representation; multi-domain fusion allows each branch to specialize. This structural complementarity explains why the single-domain CNN, despite being trained with the same physics-constrained loss, cannot match PMDF-Net's performance---it lacks the architectural vocabulary to separately reason about temporal onset and spectral deviation.

The physics-constrained loss is not merely a regularizer but an essential inductive bias that prevents the network from exploiting the SNR metric as a proxy for signal distortion. Without the feature-preservation term, an MSE-only objective drives the network toward overly smooth outputs that minimize pixel-wise error at the cost of erasing the diagnostic signatures that clinicians and inspectors rely on. The ablation study (Section~\ref{sec:results-ablation}) confirms this: removing the physics constraint degrades both SNR and feature metrics, with the envelope CCC dropping below the $0.95$ clinical relevance threshold. The physics loss is therefore not optional tuning but a necessary condition for utility in downstream inspection tasks.

From a practical deployment standpoint, the combination of high SNR improvement and minimal feature distortion enables a dramatic reduction in acquisition time. Achieving comparable defect detection accuracy with conventional signal averaging would require 256 averages; PMDF-Net delivers statistically equivalent detection from 16 averages, representing a $16\times$ reduction in laser shots. This compression of acquisition time directly translates to increased throughput in industrial inspection and reduced patient exposure in medical applications, making the proposed method not only a signal-processing advance but an operational one.